\begin{solapartat}
\punts{0,5} L'angle límit (o crític) és l'angle d'incidència respecte de la normal de la interfície pel qual l'angle de refracció és de $90\,^{\mathrm{o}}$, també respecte de la normal. A partir de la llei de Snell: $n_{bloc}\sin\theta_i = n_{aire}\sin\theta_r$.

Si l'angle refractat és $90^{\mathrm{o}}$, l'angle d'incidència compleix: $\sin\theta_{\textit{límit}} = \frac{n_{aire}}{n_{bloc}}\sin 90^{\mathrm{o}} = \frac{n_{aire}}{n_{bloc}}$.

Per tant, l'angle límit serà: $\theta_{\textit{límit}} = \arcsin\left(\frac{n_{aire}}{n_{bloc}}\right)$.

\punts{0,25} L'angle límit per a les dades del problema és: $\theta_{\textit{límit}} = \arcsin\left(\frac{1}{1,58}\right) = 39,3\,^{\mathrm{o}}$.

\punts{0,5} El fenomen de la reflexió total i l'angle límit no es pot donar si invertim els medis. En aquest cas, la llum passaria d'un índex de refracció petit a un índex gran i, per tant, l'angle refractat seria menor que l'incident. Matemàticament, es pot deduir de la llei de Snell: $n_1\sin\theta_1 = n_2\sin\theta_2$, recordant que la funció sinus creix amb l'angle per angles petits. Per al nostre cas concret, l'angle límit hauria de complir $\theta_{\textit{límit}} = \arcsin\left(\frac{1,58}{1}\right)$, que no té solució.
\end{solapartat}

\begin{solapartat}
\punts{0,75} Per poder veure la moneda, la distància ha de ser menor que la que correspon a l'angle límit. Per a distàncies més grans, tindrem reflexió total i no podrem veure la moneda.

La relació entre la distància $d$ i l'angle d'incidència $\theta_i$ és:
\[ \mathrm{tg}(\theta_i) = \frac{d}{2} \]
Calculem la distància màxima, la $d$ per l'angle límit $\theta_{\textit{límit}} = 39,3^{\mathrm{o}}$: $d_{max} = 2\cdot\mathrm{tg}(\theta_{\textit{límit}}) = 2\cdot\mathrm{tg}(39,3^{\mathrm{o}}) = 1,635\ \mathrm{m}$

\figura[0.4]{sol_fig1.png}

\punts{0,5} Esquema: angle reflectit igual a l'incident i angle refractat de $90\,^{\mathrm{o}}$

També és correcte indicar que no hi ha raig refractat.

\figura[0.4]{sol_fig2.png}
\end{solapartat}
